% Detailed provider lemma: quantitative estimates are confined below LTAR.

\subsection{Explicit grade overlap below the response boundary}
\label{subsec:grade-overlap}

We give one concrete verification of the grade-finality assertion used in
Section~\ref{sec:ltar-provider}. None of the quantities introduced here is
part of the statement of BinaryTail or of the late-source targeted response.

Let $p$ be a sufficiently large prime current after obstacle deletion and the
one global quota packing. The retained prime row has distinct simple values
and cardinality
\[
r_p\asymp \frac{p}{\log p}.
\tag{G1}
\]
Its supported grade set is
\[
I_p
=\{0\le h\le r_p:h(r_p-h)+1\ge p\}.
\]
Write
\[
I_p=[a_p,r_p-a_p]\cap\N.
\]

\begin{lemma}[Prime local interval]\label{lem:prime-local-interval}
For all sufficiently large primes,
\[
a_p=O(\log p),
\qquad
|I_p|\asymp\frac p{\log p}.
\]
\end{lemma}

\begin{proof}
Set
\[
h_p=\left\lceil\frac{2p}{r_p}\right\rceil.
\]
By (G1), $h_p=O(\log p)$ and $h_p\le r_p/2$ for large $p$. Hence
\[
h_p(r_p-h_p)\ge h_pr_p/2\ge p,
\]
so $h_p\in I_p$ and $a_p\le h_p=O(\log p)$. Symmetry of the defining
quadratic inequality gives
\[
|I_p|=r_p-2a_p+1\asymp p/\log p.
\]
\end{proof}

For a square current $Q=p^2$, the packed simple-value image has cardinality
\[
r_{p^2}\gg p/\log p.
\]
The same test shows that its D-threshold interval is nonempty and has lower
endpoint
\[
a_{p^2}=O(\log p).
\tag{G2}
\]
For $Q=p^e$, $e\ge3$, the resolver construction uses the single grade
\[
p-1.
\tag{G3}
\]

Fix once and for all a dilation $\Lambda>1$ for which the comparable-prime
estimate holds. At an endpoint $X$, let
\[
J_X=[A_X,B_X]\cap\N
\]
be the Minkowski sum of all local supported-grade sets between the fixed source
and $X$. Since sums of integer intervals are intervals, and adding singleton
resolver grades merely translates them, $J_X$ is an integer interval whenever
all local supports are nonempty.

\begin{lemma}[Prime-band width]\label{lem:prime-band-width}
There is $c>0$ such that for all sufficiently large $Y$ the prime currents in
\[
Y<p\le\Lambda Y
\]
contribute at least
\[
c\frac{Y^2}{(\log Y)^2}
\]
to the increase of interval width between endpoints $Y$ and $\Lambda Y$.
Consequently, after changing the constant,
\[
B_Y-A_Y\gg \frac{Y^2}{(\log Y)^2}
\tag{G4}
\]
for all sufficiently large $Y$.
\end{lemma}

\begin{proof}
The band $(Y,\Lambda Y]$ contains
\[
\gg Y/\log Y
\]
primes. For every such prime, Lemma~\ref{lem:prime-local-interval} gives local
width
\[
\gg Y/\log Y.
\]
Widths add under Minkowski addition, proving the first assertion. Apply the
same assertion with $Y/\Lambda$ in place of $Y$. The resulting band
$(Y/\Lambda,Y]$ is already contained in the endpoint-$Y$ construction and has
total width
\[
\gg \frac{(Y/\Lambda)^2}{(\log(Y/\Lambda))^2}
\asymp \frac{Y^2}{(\log Y)^2},
\]
which proves (G4).
\end{proof}

We next bound the lower endpoint introduced by the currents added when the
endpoint moves from $Y$ to $\Lambda Y$. It is harmless to bound this increment
by the contribution of \emph{all} currents of size at most $\Lambda Y$.

For primes $p\le\Lambda Y$, Lemma~\ref{lem:prime-local-interval} gives a lower
contribution $O(\log p)$. The crude estimate that there are at most
$\Lambda Y$ such integers gives
\[
\sum_{p\le\Lambda Y} a_p=O(Y\log Y).
\tag{G5}
\]
For square currents $p^2\le\Lambda Y$, (G2) gives
\[
\sum_{p\le\sqrt{\Lambda Y}} a_{p^2}
=O(\sqrt Y\log Y).
\tag{G6}
\]
For higher prime powers, use (G3). For every $e\ge3$, $p^e\le\Lambda Y$
forces $p\le(\Lambda Y)^{1/e}$, whence
\[
\sum_{e\ge3}\sum_{p^e\le\Lambda Y}(p-1)
=O(Y^{2/3}\log Y).
\tag{G7}
\]
The deliberately crude logarithmic factor absorbs the finite range of relevant
exponents.

Combining (G5)--(G7),
\[
A_{\Lambda Y}-A_Y=O(Y\log Y).
\tag{G8}
\]

\begin{proposition}[Eventual overlap]\label{prop:eventual-overlap}
For all sufficiently large $Y$,
\[
A_{\Lambda Y}\le B_Y+1.
\]
Consequently the union of the endpoint grade intervals contains a final
segment of $\N$.
\end{proposition}

\begin{proof}
By (G4),
\[
B_Y-A_Y
\gg
\frac{Y^2}{(\log Y)^2},
\]
whereas (G8) gives
\[
A_{\Lambda Y}-A_Y=O(Y\log Y).
\]
Since
\[
Y\log Y=o\!\left(\frac{Y^2}{(\log Y)^2}\right),
\]
we eventually have
\[
A_{\Lambda Y}-A_Y\le B_Y-A_Y+1,
\]
which is the stated overlap inequality. Thus successive endpoint intervals
overlap or are adjacent. Their lower endpoints tend to infinity because every
new nontrivial current contributes positive minimum grade. Iterating the
fixed-ratio endpoint map therefore yields a final ray.
\end{proof}

\subsection{Explicit resource summability}
\label{subsec:resource-sum}

The exact same construction gives a summable branchwise mass majorant.

At a prime current $Q=p$, a local state uses at most $r_p$ atoms, and every
retained prime atom has mass $O(p^{-2})$. Hence the local resource is
\[
O\!\left(\frac1{p\log p}\right).
\tag{R1}
\]

At a square current $Q=p^2$, every relevant local grade is at most $p-1$, while
one generic atom has mass $O(\log p/p^4)$. Hence
\[
O\!\left(\frac{\log p}{p^3}\right).
\tag{R2}
\]

At $Q=p^e$, $e\ge3$, the own/donor resolver uses $O(p)$ atoms from the current
and predecessor roles. The historical signed-inverse bounds give the
convenient majorant
\[
O\!\left(\frac{e\log p}{p^{2e-3}}\right).
\tag{R3}
\]

\begin{lemma}[Prime resource series]\label{lem:prime-resource-series}
\[
\sum_p\frac1{p\log p}<\infty.
\]
\end{lemma}

\begin{proof}
Let $X_j=\Lambda^j$. The upper comparable-prime estimate gives
\[
\#\{p:X_j<p\le X_{j+1}\}
=O\!\left(\frac{X_j}{\log X_j}\right).
\]
For $p$ in this band,
\[
\frac1{p\log p}
=O\!\left(\frac1{X_j\,j}\right).
\]
Thus the whole band contributes $O(j^{-2})$, and summing over $j$ converges.
\end{proof}

The square series
\[
\sum_p\frac{\log p}{p^3}
\]
is dominated by the corresponding integer series. For higher powers,
\[
\sum_p\sum_{e\ge3}
\frac{e\log p}{p^{2e-3}}
\]
converges because for each $p$ the inner sum is geometrically dominated and the
resulting prime sum is again bounded by a convergent integer series.

Any boundary-predecessor donor contribution is a subseries of the same
convergent majorant, and its current indices escape every fixed finite set as
the source advances.

\begin{proposition}[Vanishing Tail resource]\label{prop:resource-tail}
Let $\mathcal R(B)$ denote the total resource majorant from all currents beyond
a source base $B$. Then
\[
\mathcal R(B)\longrightarrow0
\qquad(B\to\infty).
\]
Consequently, for every $\varepsilon>0$, every sufficiently late source admits
the full grade-final BinaryTail construction with resource at most
$\varepsilon$.
\end{proposition}

\begin{proof}
Equations (R1)--(R3) define a convergent nonnegative series over the current
filtration. Tails of a convergent nonnegative series tend to zero. The grade
construction of Proposition~\ref{prop:eventual-overlap} uses only subfamilies
of the same role reservoirs, so its branchwise resource is bounded by this
common tail majorant.
\end{proof}

The proof of BinaryTail therefore separates cleanly into two consumer
statements,
\[
\text{grade image is final},
\qquad
\text{resource tail tends to zero},
\]
which are imposed on the same exact-tuple witness before physical composition.
They must not be proved using unrelated witness choices.
